\documentclass[10pt,letterpaper]{article}
\usepackage[margin=0.7in,headheight=14pt]{geometry}
\usepackage[T1]{fontenc}
\usepackage{lmodern,amsmath,amssymb,booktabs,xcolor,fancyhdr,microtype}
\usepackage[colorlinks=true,urlcolor=teal]{hyperref}
\definecolor{navy}{HTML}{192B4D}
\definecolor{teal}{HTML}{007D80}
\pagestyle{fancy}\fancyhf{}
\fancyhead[L]{\small MF 731 | Homework 3}
\fancyhead[R]{\small 05: Portfolio VaR}
\fancyfoot[C]{\small\thepage}
\setlength{\parindent}{0pt}
\setlength{\parskip}{4pt}
\newcommand{\VaR}{\operatorname{VaR}}
\hypersetup{pdftitle={MF 731 - Three-Stock Portfolio VaR}}
\begin{document}
{\Large\bfseries\textcolor{navy}{VaR of a Three-Stock Portfolio}}

\textbf{1. Data.} I use Yahoo Finance daily, unadjusted closing prices from
May 23, 2024 to May 22, 2026, inclusive. The dates match across all three
stocks: \textbf{501 prices and 500 log returns per stock}. The first and last
five observations are below (prices in dollars).

\begin{center}
\begin{tabular}{lrrr}\toprule
Date & MSFT & AAPL & GOOGL\\\midrule
2024-05-23 & 427.00 & 186.88 & 173.55\\
2024-05-24 & 430.16 & 189.98 & 174.99\\
2024-05-28 & 430.32 & 189.99 & 176.40\\
2024-05-29 & 429.17 & 190.29 & 175.90\\
2024-05-30 & 414.67 & 191.29 & 172.11\\\midrule
2026-05-18 & 423.54 & 297.84 & 396.94\\
2026-05-19 & 417.42 & 298.97 & 387.66\\
2026-05-20 & 421.06 & 302.25 & 388.91\\
2026-05-21 & 419.09 & 304.99 & 387.66\\
2026-05-22 & 418.57 & 308.82 & 382.97\\\bottomrule
\end{tabular}
\end{center}

\textbf{2. Weights and losses.} Using the market capitalizations supplied
in the question, in the order MSFT, AAPL, GOOGL,
\[
\omega=\frac{(3.109,4.536,4.640)^\top}{12.285}
=(0.253073,0.369231,0.377696)^\top.
\]
For $V=\$1{,}000{,}000$ and the daily log-return vector $x$,
\[
\ell(x)=-V\sum_i\omega_i(e^{x_i}-1),\qquad
\ell_{\rm lin}(x)=-V\omega^\top x,\qquad
\ell_{\rm quad}(x)=-V\sum_i\omega_i\left(x_i+\tfrac12x_i^2\right).
\]
For each historical loss sample, I take the smallest loss whose empirical
CDF reaches $0.95$ (the 475th sorted loss out of 500).

\textbf{3. EWMA and normal model.} Starting with $\mu_0=0$ and $\Sigma_0=0$,
I process the returns chronologically using
\[
\mu_{k+1}=0.97\mu_k+0.03x_k,\qquad
\Sigma_{k+1}=0.97\Sigma_k+0.03(x_k-\mu_k)(x_k-\mu_k)^\top.
\]
The estimates at the May 22 close are
\[
\widehat\mu=10^{-3}
\begin{pmatrix}0.472568\\3.636476\\3.305902\end{pmatrix},\qquad
\widehat\Sigma=10^{-4}
\begin{pmatrix}
3.142657&0.343139&0.034633\\
0.343139&1.926842&0.813758\\
0.034633&0.813758&4.593809
\end{pmatrix}.
\]
With $X\sim N(\widehat\mu,\widehat\Sigma)$, the linear loss is normal, so
\[
\VaR_{0.95}(\ell_{\rm lin})
=-V\omega^\top\widehat\mu
+V\sqrt{\omega^\top\widehat\Sigma\omega}\,\Phi^{-1}(0.95).
\]
For the full and quadratic losses, I use the same 100,000 normal vectors,
with random seed 731, and take their empirical 95th percentiles.

\begin{center}
\begin{tabular}{lrr}\toprule
Loss operator & Historical VaR (\$) & Normal/EWMA VaR (\$)\\\midrule
Full & 21,043.05 & 16,609.78\\
Linear & 21,415.73 & 16,869.11\\
Quadratic & 21,038.61 & 16,605.55\\\bottomrule
\end{tabular}
\end{center}

\textbf{4. Comparison.} Full and quadratic VaR differ by less than \$5
under either model. Linear VaR is about 1.6--1.8\% higher because it omits
the negative quadratic correction to loss. The larger difference is between
models: normal/EWMA full VaR is about 21\% below historical full VaR.
EWMA emphasizes recent returns, while the historical method weights all
500 returns equally; the normal assumption also changes the tail shape.

{\small These are one-trading-period forecasts. The handout labels May 25
as the next trading day, but it is Memorial Day; the next session is May 26
(\href{https://www.nyse.com/markets/hours-calendars}{NYSE calendar}). No extra
time scaling is applied. Data: \href{https://finance.yahoo.com/quote/MSFT/history/}{MSFT},
\href{https://finance.yahoo.com/quote/AAPL/history/}{AAPL},
\href{https://finance.yahoo.com/quote/GOOGL/history/}{GOOGL} on Yahoo Finance.}

\textbf{Python code and data:}
\url{https://github.com/jesuswalking/MF-731-HW1/tree/main/HW3}
\end{document}
